\section{What we learned}
\label{sec:learnings}

\begin{enumerate}[leftmargin=*,itemsep=2pt]
\item \textbf{The paper's central claims reproduce.}  The prototype neuron's activation surface,
the bit-count histogram and its convergence, the Iris overfitting contrast, perfect Flappy play
(on our clone) and the ECG denoising numbers (within 2~dB on a different signal) all come out of
a clean-room implementation.  What does not: the weight picture of Fig.~9 --- scored over 20
seeds against a criterion fixed in advance, the prototype layer does cover more of the data than
the classic one, but an untrained prototype layer already does, so the effect is the
architecture's rather than training's, and the trained angles do not decode to data patterns
(\cref{sec:iris}) --- and --- once the paper's own ECG recording reached us and removed
the signal as an excuse --- the upper half of Table~I, where our two- and three-neuron prototype
cells fall $2.5$--$3.2$~dB short ($2.2$--$3.8$~dB on the full recording) while the one-neuron cell reproduces with a \emph{larger}
prototype advantage than the paper's (\cref{sec:ecg}).  The iterative baseline of Fig.~6 was a second such failure until
we noticed that we had been grading it with the wrong quantity: a scheme that trains one input at
a time moves the per-input activations, not the superposition distribution of
\cref{eq:model-distribution}, and because the prototype block leaves the input register entangled
with the hidden qubits the two are genuinely different.  In its own observable the baseline
reaches the paper's $\approx 10^{-3}$ on four seeds of six (\cref{sec:bitcount}); only the unit
of the paper's ``iteration'' is still assumed.  Grade a scheme in the observable it optimises.

\item \textbf{Reading figures is part of reproduction.}  The bit-count target
($(\text{ones}+1)/112$), the Iris label weighting ($(\ell+1)^2$) and the angle encoding
($\theta = 2\arcsin\sqrt s$, from a printed gate angle in the companion paper) all had to be
read off figures; each was then confirmed to the figure's precision.  Where a figure was
over-read --- the Iris split ratio ``justified'' by bar counts --- an independent review caught
it and the justification was withdrawn (\cref{sec:iris}).

\item \textbf{The prototype network's advantage under noise is inherited, not intrinsic.}
Prototype and classic networks have the same two-qubit gate count and the noise displaces both
by the same amount; the prototype wins under noise because it fits better without noise.

\item \textbf{Gate noise at published levels is not the obstacle for these networks; shots
are.}  Every 7--11-qubit network survives every vendor's calibrated noise model with a
qualitatively intact result.  The Flappy controller needs $10^4$ shots per frame with a fixed
budget; the bit-count loss is biased by $31.8/N$.  The one remedy that pays now is a sequential
measurement rule ($4.6\times$ fewer shots, no retraining); margin-aware training does not work
for a 17-parameter network; common random numbers reduce the optimiser's noise by $1/c$ but
leave the estimator's bias, which is the likely (untested) limit of sampled training.

\item \textbf{Compiled circuits are $5$--$15\times$ more noise-sensitive than logical ones.}
Two native gates per controlled rotation, resynthesised single-qubit gates and routing on a
lattice all cost; a complete-bipartite network is the worst case for a lattice, and the vendor's
own compiler must be used for any superconducting submission.  Trapped ions are the natural home
of this architecture; their price is circuit time (milliseconds per shot).

\item \textbf{An optimiser that keeps a stale incumbent turns a harmless change into a
negative result.}  The ``rotating seeds break training'' finding of the first session was an
artefact of comparing losses measured on different episode batches; re-scoring the incumbent
restores $25/25$ perfect games.  Any non-stationary objective (mini-batches, hardware drift) needs
this.

\item \textbf{Penalties must be priced in the currency of the reward.}  A margin hinge that
could outweigh the per-frame survival reward produced a confidently dead controller; the
per-frame form below the per-frame reward cannot.  The same arithmetic applies to any shaping
term.

\item \textbf{Provenance has to be mechanical.}  The first hardware plan contained a dozen
wrong numbers --- a state-preparation gate count off by a factor of two, a sentence with its
inequality reversed, costs computed by hand.  The fix was a script that emits every number and
refuses to run if any calibration input lacks a source; the report now quotes only that script.
\item \textbf{A read-out calibrated on the evaluated set is transductive.}  The Iris labels are
rounded after one overall scale is fitted; fitted on the test set itself (no test labels are
used), it makes a flower's label depend on which other flowers are evaluated with it: removing
the 13 test flowers that carry three ambiguous codes changed the labels of most of the rest
($94\,\% \to 39\,\%$ on one seed).  Only the 20-seed sweep exposed it; one seed would not have.
Calibrate on the training data, state the pass criterion before running the seeds --- and when a
score is a $z$-score against each model's own baseline, do not read it as a comparison between
the models (a review caught us doing exactly that).
\end{enumerate}

\paragraph{Not done, and why.}
Noise-aware training of the 11-qubit bit-count network (13 hours for its grid) and the full
ECG/Iris noise grids (4--8 hours) exceeded the session's compute ceiling and were not run;
their per-cell costs are recorded so that a future session can decide.  Items 4--6 of the
shot-noise programme --- a loss that uses all shots, bias-corrected estimators, and amplitude
estimation --- are designed but not implemented; item 4 is the natural next step, since the
estimator bias is what now limits sampled training.  Nothing has been submitted to a quantum
computer; the plan of \cref{sec:recommendation} is ready to be.
